\documentclass[11pt]{article}
\usepackage{shared/luxcover}
\usepackage{amsmath, amssymb, amsthm}
\usepackage{mathtools}
\usepackage{bm}
\usepackage{graphicx}
\usepackage{booktabs}
\usepackage{hyperref}
\usepackage{enumitem}
\usepackage{algorithm}
\usepackage{algpseudocode}
\usepackage{xcolor}
\usepackage{multicol}
\usepackage{listings}
\hypersetup{colorlinks=true,linkcolor=black,citecolor=blue,urlcolor=blue}

\lstset{
  basicstyle=\ttfamily\scriptsize,
  breaklines=true,
  frame=single,
  numbers=left,
  numberstyle=\tiny\color{gray},
  keywordstyle=\color{blue!70!black}\bfseries,
  commentstyle=\color{green!50!black}\itshape,
  stringstyle=\color{red!60!black},
  showstringspaces=false,
}

\newtheorem{theorem}{Theorem}
\newtheorem{lemma}[theorem]{Lemma}
\newtheorem{proposition}[theorem]{Proposition}
\newtheorem{corollary}[theorem]{Corollary}
\newtheorem{definition}{Definition}
\newtheorem{property}{Property}

\title{Lux Teleport: Sovereign Omnichain Liquidity\\via Threshold Cryptography and\\Per-Chain Governance}
\author{
Zach Kelling\thanks{Corresponding author: zach@lux.network}
\quad
Vishnu Seesahai\\
\textit{Lux Industries, Inc.}\\
\texttt{research@lux.network}
}
\date{2022 (Revised December 2025)}

\begin{document}
\luxcoverpage

\begin{abstract}
We present \textbf{Lux Teleport}, a non-custodial omnichain bridge and liquidity protocol that natively connects 18 execution environments and 270 registered chain IDs through a unified MPC threshold signature network. Unlike existing bridges that rely on trusted intermediaries, multisigs, or chain-specific messaging layers, Teleport achieves trustless cross-chain asset transfer through FROST~\cite{frost} and CGGMP21~\cite{cggmp21} threshold cryptography with MPC group key verification (a single aggregate public key verified via \texttt{ecrecover} on EVM chains), per-chain nonce tracking, and automatic undercollateralization detection.

Teleport introduces three innovations: (1)~\emph{sovereign chain governance} where each connected chain's native token holders control their own bridge fees, yield strategies, and risk parameters; (2)~a \emph{yield-bearing bridge token} system where locked assets on source chains are deployed to yield strategies and the yield is reflected in bridge token share prices on destination chains; and (3)~a \emph{Shariah-compliant mode} that programmatically filters yield sources to enable the first protocol-level Islamic finance DeFi infrastructure.

The protocol supports 18 native bridge programs across every major smart contract runtime: EVM (Solidity), Solana (Anchor/Rust), TON (FunC), Sui and Aptos (Move), Cosmos/IBC (CosmWasm), Bitcoin L1 (FROST Taproot), OP\_NET (AssemblyScript), XRPL (Hooks), Polkadot (ink!), NEAR (Rust/WASM), Stellar (Soroban), Stacks (Clarity), Cardano (Aiken/Plutus), StarkNet (Cairo), ICP (Rust Canisters), Algorand (PyTeal), Fuel (Sway), and Tezos (CameLIGO). All bridge tokens implement the vendor-neutral \texttt{IBridgeToken} interface (not tied to any specific bridge vendor). The OmnichainRouter contract is non-upgradeable, with a hard-coded 1\% maximum bridge fee, configurable rotation delay (default 24h, minimum 0 for instant rotation, maximum 7 days; operational parameter, not a security invariant), and automatic pause on undercollateralization below 98.5\%. The \texttt{batchMintDeposit} function enables GPU-optimized batch processing of multiple deposits in a single transaction.

This paper supersedes Lux Teleport Protocol (2022)~\cite{luxtp2022} and Lux Bridge (2023)~\cite{luxbridge2023}.
\end{abstract}

\tableofcontents

%% ═══════════════════════════════════════════════════════════════════════════
\section{Introduction}
\label{sec:intro}
%% ═══════════════════════════════════════════════════════════════════════════

Cross-chain interoperability remains the critical unsolved problem in blockchain infrastructure. The total value locked across bridges exceeds \$15B, yet bridge exploits account for over \$2.8B in losses since 2020~\cite{rekt}. We identify four classes of existing solutions and their fundamental limitations:

\begin{enumerate}[nosep]
  \item \textbf{Trusted multisig}: Custodial risk, single points of failure (WBTC, most centralized bridges). Failure mode: compromise $m$-of-$n$ keys $\to$ drain all funds.
  \item \textbf{Optimistic}: Long withdrawal delays (7--14 days), fraud proof windows (Across, Hop). Failure mode: censored fraud proofs $\to$ invalid state accepted.
  \item \textbf{Light client}: Chain-specific implementation per pair, high complexity ($O(n^2)$ implementations for $n$ chains). Failure mode: consensus bug in light client.
  \item \textbf{MPC threshold}: Distributed key management, no single custodian, but key management complexity. Failure mode: $\geq t$ node compromise.
\end{enumerate}

Teleport addresses the MPC approach's key management complexity through two complementary threshold signature protocols:

\begin{itemize}[nosep]
  \item \textbf{FROST}~\cite{frost}: Flexible Round-Optimized Schnorr Threshold for Ed25519 and BIP-340 Taproot, providing 64-byte signatures in 2~rounds. Covers 14 chain families.
  \item \textbf{CGGMP21}~\cite{cggmp21}: Threshold ECDSA with identifiable aborts for secp256k1. Covers 6 chain families. Standard \texttt{ecrecover} on all EVM chains.
\end{itemize}

\paragraph{Contributions.} This paper makes the following contributions:
\begin{enumerate}[nosep]
  \item Formal description of the OmnichainRouter contract with provable security properties (\S\ref{sec:protocol}).
  \item Sovereign governance model separating security, economic, and protocol authority domains (\S\ref{sec:governance}).
  \item Yield-bearing bridge token architecture with 29~deployed yield strategies across 9~categories (\S\ref{sec:yield}).
  \item First protocol-level Shariah compliance framework for DeFi (\S\ref{sec:shariah}).
  \item 18 native bridge programs across 270 registered chain IDs (\S\ref{sec:chains}).
  \item Formal security analysis proving conservation, uniqueness, liveness, safety, governance isolation, and exit guarantee properties (\S\ref{sec:security}).
\end{enumerate}

%% ═══════════════════════════════════════════════════════════════════════════
\section{System Model}
\label{sec:model}
%% ═══════════════════════════════════════════════════════════════════════════

\subsection{Entities}

\begin{definition}[MPC Signer Set]
A set $\mathcal{S} = \{s_1, s_2, \ldots, s_n\}$ of $n$ nodes, each holding a share of a distributed secret key generated through distributed key generation (DKG) with no trusted dealer. A threshold $t \leq n$ defines the minimum number of honest nodes required to produce a valid signature. Default: $n = 3, t = 2$. On-chain verification uses a single \emph{MPC group address}---the aggregate public key derived from the DKG. Individual signer identities exist for accountability and identifiable abort but are not exposed on-chain: $\texttt{ecrecover}(\text{digest}, \sigma) = \texttt{mpcGroupAddress}$.
\end{definition}

\begin{definition}[OmnichainRouter]
An immutable (non-upgradeable, non-proxied) smart contract deployed on each connected chain that:
\begin{enumerate}[nosep]
  \item Verifies MPC threshold signatures for mint/release operations
  \item Maintains per-source-chain nonce bitmaps for replay prevention
  \item Distributes bridge fees to sovereign stakeholders
  \item Auto-pauses when backing ratio falls below 98.5\%
\end{enumerate}
\end{definition}

\begin{definition}[Sovereign Governor]
The native chain's DAO timelock address $G_C$ for chain $C$. Controls economic parameters:
\begin{itemize}[nosep]
  \item Bridge fee rate: $0 \leq f_b \leq 100$ bps (hard-capped, immutable maximum)
  \item Stakeholder share: $0 \leq f_s \leq 10{,}000$ bps
  \item Token registration and daily mint limits
  \item Shariah compliance mode activation
\end{itemize}
$G_C$ \emph{cannot}: mint tokens, modify nonce bitmaps, bypass collateralization checks, or shorten timelock durations.
\end{definition}

\subsection{Threat Model}

We assume a computationally bounded adversary $\mathcal{A}$ who may compromise up to $t - 1$ MPC signer nodes, control the governor address on any single chain, observe and reorder transactions on any chain, and deny service to individual relay nodes.

We assume $\mathcal{A}$ \emph{cannot} simultaneously compromise $\geq t$ MPC signer nodes, break the Discrete Logarithm Problem (DLP) or the Elliptic Curve Discrete Logarithm Problem (ECDLP), or rewrite confirmed blocks on any connected chain (finality assumption).

%% ═══════════════════════════════════════════════════════════════════════════
\section{Protocol Description}
\label{sec:protocol}
%% ═══════════════════════════════════════════════════════════════════════════

\subsection{Bridge Operations}

The protocol supports four atomic operations, each protected by threshold signature verification and nonce uniqueness.

\paragraph{Lock.} User deposits asset $A$ of amount $a$ on source chain $C_s$ into the YieldBridgeVault. The vault emits:
\[
  \texttt{Lock}(C_s, C_d, \nu, A, a, r)
\]
where $C_d$ is the destination chain, $\nu$ is a monotonically increasing nonce, and $r$ is the recipient address on $C_d$.

\paragraph{Mint.} MPC signer set $\mathcal{S}$ observes the Lock event, reaches quorum ($\geq t$ signers agree), and produces threshold signature $\sigma$ over:
\begin{equation}
  \label{eq:deposit-msg}
  m = H\bigl(\texttt{"DEPOSIT"} \,\|\, C_s \,\|\, \nu \,\|\, T \,\|\, r \,\|\, a\bigr)
\end{equation}
where $H$ is keccak256 for EVM chains or SHA-256 for non-EVM, and $T$ is the bridge token address. Any relayer submits $(C_s, \nu, T, r, a, \sigma)$ to the OmnichainRouter on $C_d$.

The router:
\begin{enumerate}[nosep]
  \item Verifies $\sigma$ against the MPC public key
  \item Checks $\mathcal{N}(C_s, \nu) = 0$ (nonce not processed)
  \item Checks daily mint limit: $\texttt{dailyMinted}[T] + a \leq \texttt{dailyMintLimit}[T]$
  \item Sets $\mathcal{N}(C_s, \nu) \leftarrow 1$ (marks nonce processed)
  \item Computes fee: $f = a \cdot f_b / 10{,}000$
  \item Mints $(a - f)$ tokens to recipient $r$
  \item Distributes fee: $f \cdot f_s / 10{,}000$ to stakeholder vault, remainder to treasury
\end{enumerate}

\paragraph{Burn.} User burns bridge tokens on $C_d$, emitting:
\[
  \texttt{Burn}(C_d, C_s, \nu', T, a, r')
\]
This is fully permissionless---no MPC signature required.

\paragraph{Release.} MPC produces $\sigma'$ over:
\begin{equation}
  m' = H\bigl(\texttt{"RELEASE"} \,\|\, C_d \,\|\, \nu' \,\|\, T \,\|\, r' \,\|\, a\bigr)
\end{equation}
Note the distinct prefix (\texttt{"RELEASE"} vs \texttt{"DEPOSIT"}) for domain separation, preventing cross-operation signature reuse.

\subsection{Nonce Management}

Each OmnichainRouter maintains a per-source-chain nonce bitmap:
\begin{equation}
  \mathcal{N}: \mathbb{Z}_{\text{chainId}} \times \mathbb{Z}_{\text{nonce}} \rightarrow \{0, 1\}
\end{equation}

The check-and-set operation $\bigl(\text{assert } \mathcal{N}(C_s, \nu) = 0; \text{ set } \mathcal{N}(C_s, \nu) \leftarrow 1\bigr)$ executes atomically within a single transaction. Combined with the deterministic message hash, this ensures each deposit is processed exactly once across all destination chains.

\subsection{Backing Attestation}

Every 12 hours, the MPC signer set attests to the total backing on each source chain:
\begin{equation}
  \sigma_b = \text{Sign}\bigl(H(\texttt{"BACKING"} \,\|\, T \,\|\, B \,\|\, \tau)\bigr)
\end{equation}
where $B$ is the total backing amount and $\tau$ is a monotonically increasing timestamp. If:
\begin{equation}
  \frac{B}{\texttt{totalMinted}[T]} < 0.985
\end{equation}
the OmnichainRouter automatically sets $\texttt{paused} \leftarrow \texttt{true}$. Unpausing requires a separate MPC signature.

%% ═══════════════════════════════════════════════════════════════════════════
\section{Threshold Cryptography}
\label{sec:threshold}
%% ═══════════════════════════════════════════════════════════════════════════

\subsection{FROST Protocol}

For chains supporting Ed25519 or Schnorr (Solana, TON, Sui, Aptos, Cosmos, Bitcoin Taproot, Polkadot, NEAR, Stellar, Cardano, Algorand, Tezos), we employ FROST~\cite{frost}.

\begin{algorithm}
\caption{FROST $t$-of-$n$ Threshold Signing}
\begin{algorithmic}[1]
\Require Message $m$, participant set $P \subseteq \mathcal{S}$, $|P| \geq t$
\Ensure Threshold Schnorr signature $\sigma = (R, s)$
\State \textbf{Round~1 (Commitment):}
\ForAll{$p_i \in P$}
  \State Sample nonce pair $(d_i, e_i) \xleftarrow{\$} \mathbb{Z}_q^2$
  \State Broadcast commitment $(D_i, E_i) = (g^{d_i}, g^{e_i})$
\EndFor
\State \textbf{Round~2 (Signature):}
\State Compute binding factor: $\rho_i = H_1(i, m, \{(D_j, E_j)\}_{j \in P})$
\State Compute group commitment: $R = \prod_{i \in P} D_i \cdot E_i^{\rho_i}$
\State Compute challenge: $c = H_2(R, Y, m)$
\ForAll{$p_i \in P$}
  \State Compute partial: $z_i = d_i + e_i \rho_i + \lambda_i s_i c$
\EndFor
\State Aggregate: $s = \sum_{i \in P} z_i$
\State \Return $\sigma = (R, s)$
\end{algorithmic}
\end{algorithm}

where $\lambda_i$ are Lagrange interpolation coefficients and $s_i$ are Shamir secret shares of the group private key.

\begin{theorem}[FROST Unforgeability~\cite{frost}]
Under the Discrete Logarithm assumption in the Random Oracle Model, FROST is existentially unforgeable under chosen-message attack for any adversary corrupting at most $t - 1$ of $n$ participants.
\end{theorem}

\subsection{CGGMP21 Protocol}

For EVM chains, TRON, XRPL, and StarkNet (secp256k1 ECDSA), we employ CGGMP21~\cite{cggmp21}. The protocol produces standard $(r, s, v)$ ECDSA signatures verifiable via \texttt{ecrecover} on all EVM chains.

\begin{theorem}[CGGMP21 Security~\cite{cggmp21}]
CGGMP21 UC-realizes ideal threshold ECDSA signing with identifiable abort under the Strong RSA, DDH, and Paillier assumptions, tolerating $t - 1$ malicious corruptions.
\end{theorem}

\subsection{Taproot Integration}

For Bitcoin L1 custody, FROST operates in Taproot mode via $\texttt{frost.KeygenTaproot}()$, producing an $x$-only public key $\tilde{X}$ (32~bytes, BIP-340~\cite{bip340}). The Taproot address $\texttt{bc1p}\langle\text{Bech32m}(\tilde{X})\rangle$ serves as the bridge vault. No on-chain smart contract is required; Bitcoin Script's native Schnorr verification validates threshold signatures.

\subsection{Signature Verification by Chain Family}

\begin{table}[h]
\centering
\small
\begin{tabular}{llll}
\toprule
\textbf{Family} & \textbf{Scheme} & \textbf{Sig.} & \textbf{Gas/Cost} \\
\midrule
EVM (30+) & ECDSA & 65\,B & 3{,}000 \\
Solana & Ed25519 & 64\,B & 0 (sysvar) \\
TON & Ed25519 & 64\,B & $\sim$0.01 TON \\
Sui & Ed25519 & 64\,B & $\sim$0.001 SUI \\
Aptos & Ed25519 & 64\,B & $\sim$100 gas \\
Cosmos & Ed25519 & 64\,B & $\sim$50{,}000 gas \\
Bitcoin & Schnorr & 64\,B & 1 vbyte \\
Polkadot & Ed25519 & 64\,B & $\sim$0.01 DOT \\
NEAR & Ed25519 & 64\,B & $\sim$0.001 N \\
Stellar & Ed25519 & 64\,B & 100 stroops \\
Cardano & Ed25519 & 64\,B & $\sim$0.17 ADA \\
StarkNet & Stark ECDSA & 64\,B & $\sim$1{,}000 gas \\
Tezos & Ed25519 & 64\,B & $\sim$0.01 XTZ \\
Algorand & Ed25519 & 64\,B & 0.001 ALGO \\
\bottomrule
\end{tabular}
\caption{Signature verification across chain families. FROST produces a single 64-byte Schnorr/Ed25519 signature regardless of the $t$-of-$n$ threshold.}
\label{tab:sigverify}
\end{table}

%% ═══════════════════════════════════════════════════════════════════════════
\section{Sovereign Chain Governance}
\label{sec:governance}
%% ═══════════════════════════════════════════════════════════════════════════

\subsection{Separation of Concerns}

Teleport cleanly separates three authority domains:

\begin{table}[h]
\centering
\small
\begin{tabular}{lll}
\toprule
\textbf{Domain} & \textbf{Controller} & \textbf{Scope} \\
\midrule
Security & MPC signer set & Mint/release authorization \\
Economics & Chain governor & Fees, limits, strategies \\
Protocol & Immutable code & Nonce checks, verification \\
\bottomrule
\end{tabular}
\caption{Authority domain separation. No domain can override another.}
\label{tab:separation}
\end{table}

\begin{property}[Governance Isolation]
\label{prop:isolation}
For any governor $G_C$ of chain $C$:
\begin{enumerate}[nosep]
  \item $G_C$ cannot call \texttt{bridgeMint} without valid MPC signature $\sigma$
  \item $G_C$ cannot modify $\mathcal{N}$ (nonce bitmap) directly
  \item $G_C$ cannot set $f_b > 100$ bps (hard-coded maximum)
  \item $G_C$ cannot change stakeholder vault without 48-hour timelock
  \item $G_C$ cannot change MPC signers (MPC-controlled, configurable delay up to 7 days)
\end{enumerate}
\end{property}

\subsection{Per-Chain Deployment}

Each chain deploys an independent OmnichainRouter with its own governance:

\begin{center}
\small
\begin{tabular}{lll}
\toprule
\textbf{Chain} & \textbf{Governor} & \textbf{Fee Recipient} \\
\midrule
Lux C-Chain & LUX DAO & xLUX (LiquidLUX) \\
Zoo EVM & ZOO DAO & ZOO Stakers \\
Third-party L1 & Their DAO & Their staking vault \\
\bottomrule
\end{tabular}
\end{center}

The MPC signer set can be shared across chains or each chain can maintain its own. In either case, the fee flow is sovereign to each chain.

\subsection{Signer Rotation}

The rotation delay is a configurable operational parameter set by the governor (default 24 hours, minimum 0 for instant rotation, maximum 7 days). The delay is an operational convenience---not a security invariant---because the MPC threshold \emph{is} the trust model. A compromised governor cannot rotate signers; only the current $t$-of-$n$ MPC set can initiate rotation.

\begin{enumerate}[nosep]
  \item Current $t$-of-$n$ MPC set signs rotation proposal with new MPC group address.
  \item Configurable delay begins (default 24h, governor-settable $\in [0, 7\text{d}]$). Pending signers visible on-chain.
  \item During delay: users can exit (burn tokens $\to$ receive underlying).
  \item After delay expires: anyone calls \texttt{executeSignerRotation()}.
  \item Current signers can cancel during delay via \texttt{cancelSignerRotation()}.
\end{enumerate}

%% ═══════════════════════════════════════════════════════════════════════════
\section{Yield-Bearing Bridge Tokens}
\label{sec:yield}
%% ═══════════════════════════════════════════════════════════════════════════

\subsection{Share-Based Accounting}

Yield-bearing bridge tokens (yLETH, yLBTC, yLUSD) use ERC-4626-like~\cite{erc4626} share accounting:
\begin{equation}
  \text{sharePrice}(t) = \frac{\text{totalBacking}_{\text{source}}(t)}{\text{totalShares}_{\text{dest}}(t)}
\end{equation}

When yield accrues on the source chain, the MPC attests to new backing via Equation~\eqref{eq:deposit-msg} with \texttt{"BACKING"} prefix. The share price increases automatically---holders earn yield without transactions.

\subsection{Strategy Categories}

29 deployed strategies span 9 categories:

\begin{table}[h]
\centering
\scriptsize
\begin{tabular}{llr}
\toprule
\textbf{Category} & \textbf{Protocols} & \textbf{APY} \\
\midrule
Liquid Staking & Lido, Rocket Pool & 3--5\% \\
Restaking & EigenLayer, Symbiotic, Karak & 3--8\% \\
Lending & Aave, Compound, Morpho, Euler, Spark, Fluid & 2--12\% \\
Stablecoin & MakerDAO/Sky, Ethena, Frax & 4--15\% \\
LP/DEX & Curve, Convex, Pendle, L2 DEXs & 5--20\% \\
BTC Native & Babylon, Lombard, SolvBTC, CoreDAO & 3--12\% \\
Solana & Marinade, Jito, Kamino & 5--10\% \\
TON & Tonstakers, Bemo, STON.fi & 5--15\% \\
Perps LP & LPX fee-based LP & 10--30\% \\
\bottomrule
\end{tabular}
\caption{Yield strategy taxonomy. 29 strategies across 9 categories.}
\label{tab:strategies}
\end{table}

\subsection{Yield Composition}

Bridge token holders earn layered yield:
\begin{equation}
  \text{APY}_{\text{total}} = \underbrace{\text{APY}_{\text{base}}}_{\text{staking}} + \underbrace{\text{APY}_{\text{re}}}_{\text{restaking}} + \underbrace{\text{APY}_{\text{bridge}}}_{\text{fee share}} + \underbrace{\text{APY}_{\text{DeFi}}}_{\text{LP/lending}}
\end{equation}

For yLBTC on Lux: $\sim$5\% (Babylon) + $\sim$2\% (EigenLayer) + $\sim$1\% (xLUX bridge fees) + $\sim$5--15\% (DeFi) = 13--23\% total.

%% ═══════════════════════════════════════════════════════════════════════════
\section{Shariah Compliance}
\label{sec:shariah}
%% ═══════════════════════════════════════════════════════════════════════════

Islamic finance law (Shariah) prohibits \emph{riba} (interest/usury), affecting $\sim$1.8~billion Muslims. Teleport introduces the \texttt{ShariaFilter} contract for protocol-level compliance:

\begin{definition}[Compliance Classification]
Each yield strategy $Y$ is assigned a status $\mathcal{C}(Y) \in \{\textsc{Halal}, \textsc{Haram}, \textsc{Conditional}, \textsc{UnderReview}\}$.
\end{definition}

\begin{table}[h]
\centering
\small
\begin{tabular}{lll}
\toprule
\textbf{Status} & \textbf{Examples} & \textbf{Rationale} \\
\midrule
\textsc{Halal} & DEX fees, bridge fees & Service fee \\
& Validator staking & Network security \\
& LP provision & Liquidity service \\
& Babylon BTC staking & Fee-based security \\
\midrule
\textsc{Haram} & Aave/Compound & Riba (interest) \\
& MakerDAO DSR & Savings interest \\
& sDAI, USDY & Interest-bearing \\
\midrule
\textsc{Conditional} & Lido, EigenLayer & Structure-dependent \\
\bottomrule
\end{tabular}
\caption{Shariah compliance classification of yield sources}
\label{tab:shariah}
\end{table}

A Shariah Advisory Board (SAB) multisig updates classifications. When \texttt{shariahMode} is enabled, only $\mathcal{C}(Y) = \textsc{Halal}$ strategies receive deposits.

%% ═══════════════════════════════════════════════════════════════════════════
\section{Supported Execution Environments}
\label{sec:chains}
%% ═══════════════════════════════════════════════════════════════════════════

18 native bridge programs are deployed across the following execution environments:

\begin{table}[h]
\centering
\scriptsize
\begin{tabular}{rllll}
\toprule
\textbf{\#} & \textbf{Runtime} & \textbf{Lang.} & \textbf{Sig.} & \textbf{Chains} \\
\midrule
1 & EVM & Solidity & ECDSA & 180+ \\
2 & SVM & Rust & Ed25519 & 3 \\
3 & TVM (TON) & FunC & Ed25519 & 2 \\
4 & Sui Move & Move & Ed25519 & 3 \\
5 & Aptos Move & Move & Ed25519 & 3 \\
6 & CosmWasm & Rust & Ed25519 & 20+ \\
7 & OP\_NET & AS & Schnorr & 1 \\
8 & Bitcoin & FROST & Schnorr & 3 \\
9 & XRPL Hooks & C & ECDSA & 3 \\
10 & Substrate & Rust & Ed25519 & 14 \\
11 & NEAR & Rust & Ed25519 & 2 \\
12 & Soroban & Rust & Ed25519 & 2 \\
13 & Clarity & Clarity & secp256k1 & 2 \\
14 & Plutus & Aiken & Ed25519 & 3 \\
15 & Cairo & Cairo & Stark & 1 \\
16 & ICP & Rust & t-ECDSA & 1 \\
17 & AVM & PyTeal & Ed25519 & 2 \\
18 & FuelVM & Sway & ECDSA & 2 \\
19 & Michelson & LIGO & Ed25519 & 2 \\
20 & TVM (TRON) & Solidity & ECDSA & 3 \\
\bottomrule
\end{tabular}
\caption{18 native bridge programs, 270 registered chain IDs. All bridge tokens implement the vendor-neutral \texttt{IBridgeToken} interface.}
\label{tab:runtimes}
\end{table}

%% ═══════════════════════════════════════════════════════════════════════════
\section{Formal Security Properties}
\label{sec:security}
%% ═══════════════════════════════════════════════════════════════════════════

\begin{property}[Conservation]
For every bridge token $T$ at time $t$:
\begin{equation}
  \texttt{totalMinted}[T](t) \leq \texttt{totalBacking}[T](t)
\end{equation}
Enforced by automatic pause when $B / M < 0.985$.
\end{property}

\begin{property}[Uniqueness]
For all source chains $C_s$ and nonces $\nu$:
\begin{equation}
  \bigl|\{tx : tx \text{ processes } (C_s, \nu)\}\bigr| \leq 1
\end{equation}
Enforced by atomic check-and-set on nonce bitmap.
\end{property}

\begin{property}[Liveness]
If $\geq t$ of $n$ MPC nodes are online and at least one relayer is operational, the bridge processes new deposits within one source-chain confirmation period plus one destination-chain block time.
\end{property}

\begin{property}[Safety]
Under the DLP/ECDLP assumptions:
\begin{equation}
  |\mathcal{A} \cap \mathcal{S}| < t \implies \Pr[\text{forge}(\sigma)] \leq \text{negl}(\lambda)
\end{equation}
\end{property}

\begin{property}[Exit Guarantee]
Users can burn bridge tokens and receive underlying assets within $\max(\text{rotationDelay}, \text{confirm\_time})$ under all conditions, including governor compromise and single MPC node failure. The rotation delay is configurable (default 24h, max 7 days).
\end{property}

%% ═══════════════════════════════════════════════════════════════════════════
\section{Fee Economics}
\label{sec:economics}
%% ═══════════════════════════════════════════════════════════════════════════

\subsection{Fee Structure}

\begin{align}
  \text{mintAmount} &= \text{deposit} \cdot \left(1 - \frac{f_b}{10{,}000}\right) \\
  \text{toStakeholders} &= f_b \cdot \text{deposit} \cdot \frac{f_s}{10{,}000^2} \\
  \text{toTreasury} &= f_b \cdot \text{deposit} / 10{,}000 - \text{toStakeholders}
\end{align}

Default: $f_b = 10$ bps (0.1\%), $f_s = 9{,}000$ bps (90\% to stakeholders).

\subsection{xLUX Flywheel}

On Lux, the stakeholder vault is LiquidLUX (xLUX). All fees compound:
\begin{equation}
  \text{APY}_{\text{xLUX}} = \frac{\sum_p \text{fees}_p}{\text{totalStaked}_{\text{LUX}}}
\end{equation}
where $p$ ranges over bridge, DEX, lending, perps, and NFT AMM protocols. This creates a positive feedback loop: more bridged assets $\to$ more yield $\to$ more fees $\to$ higher xLUX APY $\to$ more LUX staked $\to$ deeper liquidity.

%% ═══════════════════════════════════════════════════════════════════════════
\section{Related Work}
\label{sec:related}
%% ═══════════════════════════════════════════════════════════════════════════

\textbf{Wormhole}~\cite{wormhole} uses a 13-of-19 guardian multisig. Teleport's 2-of-3 MPC is more capital-efficient but trades off decentralization breadth for operational simplicity and faster key rotation.

\textbf{LayerZero}~\cite{layerzero} provides generic messaging (Oracle + Relayer). Teleport is bridge-specific, optimizing for asset transfer with yield composition.

\textbf{IBC}~\cite{ibc} achieves trust-minimized Cosmos-to-Cosmos transfers via light client verification. Teleport's CosmWasm bridge integrates with IBC for inter-Cosmos transfers while using MPC for non-IBC chains.

\textbf{tBTC}~\cite{tbtc} pioneered threshold ECDSA for Bitcoin bridging. Teleport generalizes this to 18 native bridge programs across 270 chain IDs with yield-bearing tokens and sovereign governance.

\textbf{Axelar}~\cite{axelar} runs a dedicated PoS chain for cross-chain communication. Teleport avoids the overhead of a separate consensus chain by embedding MPC signing in existing infrastructure.

%% ═══════════════════════════════════════════════════════════════════════════
\section{Conclusion}
\label{sec:conclusion}
%% ═══════════════════════════════════════════════════════════════════════════

Lux Teleport provides the infrastructure for a sovereign omnichain liquidity layer: non-custodial, non-upgradeable, per-chain governed, optionally Shariah-compliant, and natively deployed across every major blockchain execution environment. By separating attestation (MPC), governance (per-chain DAO), and protocol logic (immutable code), Teleport enables any blockchain to join the liquidity network while retaining full sovereignty over its economic parameters.

The yield-bearing bridge token system makes bridged capital productive across 29 strategies, the Shariah compliance mode opens DeFi to 1.8~billion Muslims worldwide, and the 270-chain registry covers every blockchain in the top~200 by TVL.

All source code is open-source at \texttt{github.com/luxfi}.

%% ═══════════════════════════════════════════════════════════════════════════
%% References
%% ═══════════════════════════════════════════════════════════════════════════
\begin{thebibliography}{99}

\bibitem{frost} C.~Komlo and I.~Goldberg, ``FROST: Flexible Round-Optimized Schnorr Threshold Signatures,'' in \emph{Selected Areas in Cryptography (SAC)}, 2020.

\bibitem{cggmp21} R.~Canetti, R.~Gennaro, S.~Goldfeder, N.~Makriyannis, and U.~Peled, ``UC Non-Interactive, Proactive, Threshold ECDSA with Identifiable Aborts,'' in \emph{ACM CCS}, 2020.

\bibitem{bip340} P.~Wuille, J.~Nick, and T.~Ruffing, ``BIP-340: Schnorr Signatures for secp256k1,'' Bitcoin Improvement Proposal, 2020.

\bibitem{bip341} P.~Wuille, J.~Nick, and A.J.~Towns, ``BIP-341: Taproot: SegWit version~1 spending rules,'' Bitcoin Improvement Proposal, 2020.

\bibitem{eip712} R.~Floersch and J.~Baylina, ``EIP-712: Typed Structured Data Hashing and Signing,'' Ethereum Improvement Proposal, 2017.

\bibitem{erc4626} J.~Bebis \emph{et al.}, ``EIP-4626: Tokenized Vault Standard,'' Ethereum Improvement Proposal, 2022.

\bibitem{wormhole} Wormhole Foundation, ``Wormhole: A Generic Message Passing Protocol,'' Whitepaper, 2022.

\bibitem{layerzero} R.~Zarick, B.~Pellegrino, and C.~Banister, ``LayerZero: An Omnichain Interoperability Protocol,'' Whitepaper, 2022.

\bibitem{ibc} C.~Goes \emph{et al.}, ``The Interblockchain Communication Protocol,'' Cosmos Documentation, 2020.

\bibitem{tbtc} M.~Luongo and C.~Ream, ``tBTC: A Decentralized Redeemable BTC-backed ERC-20 Token,'' Keep Network, 2020.

\bibitem{axelar} S.~Layr \emph{et al.}, ``Axelar: Connecting Web3,'' Whitepaper, 2022.

\bibitem{rekt} Rekt News, ``Leaderboard,'' \url{https://rekt.news/leaderboard/}, 2024.

\bibitem{luxtp2022} Z.~Kelling and V.~Seesahai, ``Lux Teleport Protocol,'' Lux Industries, Inc., 2022. \emph{(Superseded by this paper.)}

\bibitem{luxbridge2023} Z.~Kelling and V.~Seesahai, ``Lux Bridge: Threshold-Secured Cross-Chain Communication with MPC Custody,'' Lux Industries, Inc., 2023. \emph{(Superseded by this paper.)}

\bibitem{luxthreshold2021} Z.~Kelling, ``Lux Universal Threshold Signatures,'' Lux Industries, Inc., 2021. \emph{(Companion paper: threshold crypto details.)}

\end{thebibliography}

\end{document}
